\documentclass[a4paper,11pt]{ltjsarticle}
\usepackage[margin=25mm,headheight=16pt]{geometry}
\usepackage{amsmath,amssymb,amsthm,mathtools}
\usepackage{enumitem,booktabs}
\usepackage[unicode,hidelinks]{hyperref}
\usepackage{fancyhdr}
\usepackage{microtype}
\setlength{\emergencystretch}{2em}
\allowdisplaybreaks[2]
\pagestyle{fancy}\fancyhf{}
\fancyhead[L]{\small 自己同値付きdg圏のétale局所実現}
\fancyhead[R]{\small 2026-11-13 / v2}
\fancyfoot[C]{\thepage}
\newcommand{\CC}{\mathbb C}
\newcommand{\OO}{\mathcal O}
\newcommand{\Ku}{\operatorname{Ku}}
\newcommand{\Perf}{\operatorname{Perf}}
\newcommand{\RHom}{\operatorname{RHom}}
\newcommand{\otL}{\mathbin{\otimes^{\mathbf L}}}
\newtheorem{theorem}{定理}[section]
\newtheorem{lemma}[theorem]{補題}
\newtheorem{proposition}[theorem]{命題}
\newtheorem{corollary}[theorem]{系}
\theoremstyle{definition}
\newtheorem{definition}[theorem]{定義}
\newtheorem{example}[theorem]{例}
\theoremstyle{remark}
\newtheorem{remark}[theorem]{注意}
\renewcommand{\proofname}{証明}
